\scanpage{379}
\mspage{\deleted{\unsure{117}}118}
% archive sheet number: 378
le foncteur
\[
\underline t\mathrel{\unsure{\mapstochar\relbar|\mapsto}}C_\infty(\underline t),\quad
\underline{\mathrm{Hom\,cof}}(\mathbb N^\circ,T)\longrightarrow(\mathrm{Hot})
\]
se factorise par le groupoïde fondamental
\[
\Pi\bigl(\underline{\mathrm{Hom\,cof}}(\mathbb N^\circ,T)\bigr)\mathrel{\unsure{-|\to}}(\mathrm{Hot}).
\]

Quand ce groupoïde, i.e. \inserted{quand} la catégorie $\underline{\mathrm{Hom\,cof}}$ elle-même est connexe, on peut \unsure{donc|dé-} conclure que tous les $C_\infty(\underline t)$ sont isomorphes entre eux dans $\mathrm{Hot}$. Pour conclure même à un système \uline{transitif} d'isom. entre les $C_\infty(\underline t)$, i.e. que $C_\infty(\underline t)$ est ``unique à isom. \uline{unique} près'', il faudrait savoir \inserted{de plus} que ce groupoïde fondamental est \deleted{\unsure{simplement c}} équivalent à la catégorie ponctuelle (groupe fondamental trivial), ou, comme on dit, que $\underline{\mathrm{Hom}}\ \underline{\mathrm{cof}}$ soit \uline{simplement connexe}, ou $1$-connexe.

La connexité de $\underline{\mathrm{Hom\,cof}}$ est assurée quand \inserted{la catégorie} $T$ a (comme c'est le cas pour $\unsure{\underline{\mathrm{Tg}}|\underline{T_g}}$), une loi de composition $t,t'\mathrel{\unsure{\to|-}}t*t'$, avec des morphismes fonctoriels en $t,t'$
\diagram{\begin{tikzcd}[ampersand replacement=\&,row sep=small,column sep=normal]
\&t\\[-.6em]
t*t'\arrow[ur]\arrow[dr]\&\\[-.6em]
\&t'
\end{tikzcd}}
(cette dernière condition étant \unsure{automatique}
